% RecaLL final report. Compile on Overleaf with pdfLaTeX or XeLaTeX (Menu > Compiler).
% If the page numbers or contents look wrong after the first run, recompile once.
\documentclass[12pt,a4paper]{article}

\usepackage[T1]{fontenc}
\usepackage{lmodern}
\usepackage[margin=30mm]{geometry}
\usepackage{graphicx}
\usepackage{float}
\usepackage{array}
\usepackage{enumitem}
\usepackage{xcolor}
\usepackage{fancyhdr}
\usepackage{lastpage}
\usepackage{setspace}
\usepackage{hyperref}

% Text and links are black. Change this one colour to blue to match the FocusGuard report.
\definecolor{linkcol}{HTML}{000000}
\hypersetup{colorlinks=true,linkcolor=linkcol,urlcolor=linkcol,citecolor=linkcol,
  pdftitle={RecaLL: An AI-Powered Academic Learning Companion},
  pdfauthor={Al Rubayet, Md. Rony Hossain, Imteajul Islam, Tahsin Ahmed Chowdhury}}

\setcounter{tocdepth}{3}
\setcounter{secnumdepth}{3}
\graphicspath{{figures/}}
\sloppy
\setlength{\parskip}{0pt}

\newcommand{\code}[1]{\texttt{#1}}
\newcommand{\shot}[4][\linewidth]{%
  \begin{figure}[H]
    \centering
    \includegraphics[width=#1]{#2}
    \caption{#3}
    \label{#4}
  \end{figure}}

\pagestyle{fancy}
\fancyhf{}
\renewcommand{\headrulewidth}{0pt}
\fancyfoot[R]{\thepage\ of \pageref{LastPage}}

\begin{document}

% ------------------------------------------------------------------ Cover
\pagenumbering{alph}
\begin{titlepage}
  \centering
  {\Large Software Development Project (Sessional) [CSE-300]\par}
  \vspace{10mm}
  \includegraphics[width=42mm]{cuet_logo}\par
  \vspace{14mm}
  {\Large Final Report on:\par}
  \vspace{2mm}
  {\LARGE\bfseries RecaLL: An AI-Powered Academic\\ Learning Companion\par}
  \vspace{14mm}
  \begin{tabular}{|p{70mm}|p{35mm}|}
    \hline
    \textbf{Name} & \textbf{Student ID} \\ \hline
    Al Rubayet & 2204002 \\ \hline
    Md. Rony Hossain & 2204005 \\ \hline
    Imteajul Islam & 2204015 \\ \hline
    Tahsin Ahmed Chowdhury & 2204028 \\ \hline
  \end{tabular}\par
  \vspace{14mm}
  Submitted to:\par
  \textbf{Sharmistha Chanda Tista}\par
  Assistant Professor,\par
  Department of CSE, CUET\par
  \vfill
  \rule{\linewidth}{0.6pt}\par
  \vspace{2mm}
  {\large Department of Computer Science and Engineering (CSE)\par}
  {\large\bfseries Chittagong University of Engineering \& Technology (CUET)\par}
  {\large Chattogram-4349, Bangladesh\par}
\end{titlepage}

% ------------------------------------------------------------------ Contents
\pagenumbering{roman}
\setcounter{page}{1}
\pagestyle{empty}
\tableofcontents
\clearpage
\pagestyle{fancy}
\pagenumbering{arabic}
\setcounter{page}{1}

% ------------------------------------------------------------------ 1
\section{Title}
RecaLL: An AI-Powered Academic Learning Companion

% ------------------------------------------------------------------ 2
\section{Introduction}
University students spend much of their study time inside PDF textbooks and lecture
slides. Reading is comfortable, but it gives little evidence of whether anything has
actually been understood. A chapter can feel clear on the page and still be forgotten
by the time of the examination. Good study techniques, such as structured note-taking,
self-testing and spaced revision, are well known, yet doing them by hand takes more
time than most students can give to every topic.

\textbf{RecaLL} was developed to close this gap. It is a desktop application that takes
an academic PDF and turns it into a guided study path: a table of contents and topic
tree, handwritten-style Cornell notes for each topic, Socratic questions that probe
understanding, flashcards scheduled by the Free Spaced Repetition Scheduler (FSRS),
and a chat assistant whose answers cite the uploaded documents.

This report documents the design, architecture and implementation of RecaLL as
delivered. It covers the background and motivation of the project, its layered
architecture, a walkthrough of each implemented module, a comparison with existing
tools, and the conclusions drawn from the development process.

% ------------------------------------------------------------------ 3
\section{Background and Present State}
Software that supports academic study generally falls into three categories. The first
is the \textbf{PDF reader}, which displays a document and allows highlighting but offers
no way to check understanding. The second is the \textbf{document chat tool}, which
answers questions about an uploaded file. These tools are fast and convenient, but the
student remains the one asking, and a good answer from the tool does not show that the
student could have produced it. The third is the \textbf{flashcard application}, which
is effective for long-term retention through spaced repetition but depends on the
quality of the cards, and writing good cards from a long textbook is slow.

Recent products have begun to combine these categories, and Section~\ref{sec:related}
compares several of them. What remains less common is a single tool that follows the
material from the original PDF through to questioning and scheduled revision, runs as a
desktop application, and lets the student choose between cloud and local AI models.
RecaLL was built around this complete path.

% ------------------------------------------------------------------ 4
\section{Motivation}
The motivation for RecaLL came from four recurring problems in ordinary study:

\begin{itemize}[leftmargin=*]
  \item \textbf{Passive learning.} Reading a PDF encourages recognition rather than
  recall, so a topic can seem understood without being tested.
  \item \textbf{Time-consuming note-taking.} Preparing structured notes from a textbook
  takes significant effort, particularly for mathematical and technical subjects.
  \item \textbf{Poor retention.} Without systematic revision, previously learned
  concepts fade, and students often notice this only close to an examination.
  \item \textbf{Lack of personalised assessment.} Ordinary study materials cannot
  identify the specific misunderstanding behind a wrong explanation.
\end{itemize}

A single tool that handles the mechanical work of structuring, note-making and
scheduling leaves the student more time for the part that builds understanding:
explaining ideas and being questioned on them.

% ------------------------------------------------------------------ 5
\section{Specific Objectives and Outcomes}
The specific objectives of the project were:

\begin{itemize}[leftmargin=*]
  \item To automatically extract the structure of an academic PDF and organise it into
  chapters and topics.
  \item To generate structured, handwritten-style Cornell notes for individual topics,
  with mathematical notation rendered correctly.
  \item To evaluate understanding through Socratic questions that test reasoning rather
  than simple recall.
  \item To identify knowledge gaps and misconceptions and turn them into targeted
  flashcards.
  \item To schedule flashcard revision using the FSRS algorithm.
  \item To answer questions from the uploaded documents, with citations to the sources.
  \item To provide generated quizzes and a transparent view of learning progress.
  \item To support several AI providers, including an optional local model, with
  personal API keys stored in encrypted form.
\end{itemize}

The outcome is a working desktop application that delivers these features through one
interface, as described in Section~\ref{sec:implementation}.

% ------------------------------------------------------------------ 6
\section{Outline of Methodology}
RecaLL follows a layered architecture in which each layer has a clear responsibility.
The desktop interface communicates with a local backend, the backend owns all processing
and storage, and every AI request passes through a single provider layer.
Figure~\ref{fig:architecture} shows this structure, and Figure~\ref{fig:workflow} shows
the overall workflow followed by a student.

\shot{architecture}{System architecture of RecaLL}{fig:architecture}

\shot{workflow}{System workflow}{fig:workflow}

\subsection{Desktop and User Interface Layer}
The application is packaged as a desktop program using Tauri~v2. Its interface is built
with React~19 and TypeScript, styled with Tailwind~CSS, with React Router handling
navigation and Redux Toolkit managing shared state. Mathematical notation is rendered
with KaTeX, and PDF pages are displayed with PDF.js (\code{pdfjs-dist}).

\subsection{Backend API Layer}
A FastAPI backend written in Python runs locally and serves the interface over HTTP.
It exposes two groups of routes. The application API under \code{/api/v1} covers
documents, conversations, flashcards, quizzes, reviews and learning analytics, and uses
token-based authentication. A second group of study routes covers books, topics, notes,
annotations and the Socratic drill.

\subsection{Document Processing Layer}
Uploaded documents pass through a staged pipeline: validation, extraction,
normalisation, chunking, embedding and indexing. Text extraction is implemented as a
plugin system with PyMuPDF4LLM as the default, and Marker available as an optional
extractor for scanned or handwritten documents.

\subsection{AI Provider Layer}
All providers implement a common interface, \code{BaseLLMProvider}, and a factory selects
the provider for each request, so features do not depend on any one vendor. Gemini,
OpenAI, Groq and a local Ollama server are supported. The layer also applies rate
limiting and retries transient failures with increasing delays. Prompts are kept in a
central registry, and each can be customised without changing the code.

\subsection{Storage Layer}
Data is stored locally. SQLite in Write-Ahead Logging (WAL) mode holds an application
database for accounts, documents, chunks, conversations, flashcards, quizzes and reviews,
and a per-user study database for books, topics, notes, annotations and drill records.
Uploaded PDFs and images are kept in local file storage.

\subsection{Core Features}
The features delivered on this architecture are document import and processing, a PDF
reader with topic navigation, AI-generated Cornell notes, Socratic drills, misconception
detection with flashcard generation, FSRS-based review, the Knowledge Hub, generated
quizzes, a dashboard with learning progress, and AI provider and security settings.

% ------------------------------------------------------------------ 7
\section{Implementation}
\label{sec:implementation}
This section describes each implemented module in terms of its frontend, its backend and
its resulting functionality. The screenshots were taken from the running application
using documents uploaded by the team, including a CUET lecture on Artificial Intelligence.

% ---- 7.1
\subsection{Document Import and Processing}
\subsubsection{Frontend}
The Documents page lists every uploaded document as a card showing its type, age, page
count and processing status. Students can search by title, filter by status (All, Ready,
Processing or Failed) and sort the list. The Add Knowledge action starts an upload, and a
progress dialog reports each stage of processing.

\subsubsection{Backend}
The \code{/books/upload} route reads the PDF's built-in table of contents and, when none
exists, falls back to detecting headings from the text. Front-matter entries such as
covers and copyright pages are filtered out before the topic tree is built. A document
identified as scanned is routed to the Marker extractor when it is available. The
processing pipeline records its current stage and progress in a job table, so that
progress can be reported to the interface. The chunker splits text along headings,
paragraphs and page boundaries, with a configurable overlap between chunks.

\subsubsection{Functionality}
A student uploads a PDF and receives a navigable topic tree along with text chunks that
later features use as source material.

\shot{pipeline_flow}{Document processing pipeline}{fig:pipeline}

\shot{documents}{Uploaded documents with processing status}{fig:documents}

% ---- 7.2
\subsection{PDF Reader and Topic Navigation}
\subsubsection{Frontend}
The Document Reader presents the topic tree in a sidebar with a filter box. Selecting a
topic shows its study actions: View PDF, Extract Concepts, Generate Flashcards, Study
Notes and View Markdown. The reader displays the original PDF pages and supports
highlights and annotations. An AI chat sidebar is available beside the document.

\subsubsection{Backend}
Annotations are saved per book through the \code{/books/\{book\_id\}/annotations} routes.
Further routes explain a highlighted passage or generate flashcards from it, and the
annotated document can be exported as a PDF.

\subsubsection{Functionality}
Students can read the source material and study from it in one place, moving between a
topic and its exact page, as shown in Figure~\ref{fig:reader}.

\shot{reader_pdf}{PDF reader beside the detected topics}{fig:reader}

% ---- 7.3
\subsection{AI-Generated Cornell Notes}
\subsubsection{Frontend}
The notes editor presents notes on notebook paper with ruled, grid or plain styles, a
handwritten font and highlighter effects, and renders mathematics with KaTeX. Students can
edit notes, add images and generate a Cornell note with one button, while a progress bar
reports the state of generation. A separate Notes Hub groups notes into smart views and by
book.

\subsubsection{Backend}
The \code{/topics/\{topic\_id\}/notes/scaffold-stream} route generates the note and streams
progress events to the interface. It uses the \code{cornell\_scaffold\_prompt} template,
which asks for a fixed structure: a core intuition, invariants and definitions, step-by-step
mechanisms, a worked example, visual comparisons, common exam pitfalls and self-testing cue
questions. Notes are stored against their topic.

\subsubsection{Functionality}
Each topic can be given a complete, consistently structured note in seconds, which the
student can then edit. Figures~\ref{fig:cornell} and~\ref{fig:cornell2} show a note generated
from one slide of a CUET lecture.

\shot[0.8\linewidth]{cornell}{Cornell note: core intuition and definitions}{fig:cornell}

\shot[0.8\linewidth]{cornell_analogy}{Cornell note: analogy and first-principles breakdown}{fig:cornell2}

% ---- 7.4
\subsection{Socratic Diagnostic Drills}
\subsubsection{Frontend}
The Socratic Drill widget appears on each topic and can also target one extracted concept.
It presents one question at a time, labelled with its question type, together with an
optional hint and a text box in which the student explains their reasoning.

\subsubsection{Backend}
The \code{/topics/\{topic\_id\}/drill/generate} route sends the topic content to the selected
provider with the \code{socratic\_question\_prompt} template. Each drill contains exactly two
questions drawn from three types: causal mechanism, counterfactual and applied scenario.
Simple recall questions are excluded by the prompt. Each question carries the key invariants
that a good answer should address and a leading hint. Because model output can be malformed,
the response passes through a repair step and is validated against a Pydantic schema.

\subsubsection{Functionality}
Instead of asking the student to repeat a definition, the drill asks why something works and
what would happen if conditions changed. Figure~\ref{fig:drill} shows a question generated from
the lecture on rational agents.

\shot{drill_flow}{Socratic drill flow}{fig:drillflow}

\shot{socratic_drill}{Drill question generated from a CUET lecture}{fig:drill}

% ---- 7.5
% TEAM: this section describes grading as designed. In the code reviewed for this report,
% app/socratic_drill.py requests "socratic_evaluation_prompt", which is NOT registered in
% app/prompt_manager.py, so grading receives an empty prompt. Fix this before submitting,
% or rewrite this section and the matching text in the Conclusion.
\subsection{Misconception Detection and Flashcard Generation}
\subsubsection{Frontend}
After an answer is submitted, the widget shows a mastery score, the knowledge gaps, the
misconceptions and a Socratic nudge. Students can pin the gaps to their study notes and save
the suggested flashcards to their review deck.

\subsubsection{Backend}
The \code{/topics/\{topic\_id\}/drill/evaluate} route compares the answer against the question's
key invariants and the source content. The score falls into one of four bands: mastered
(85--100), developing (60--84), fragile (35--59) and misconception (0--34). Each gap includes
an explanation and the reason it matters, and each misconception includes the pitfall, the
correct theory and an examination tip. The number of suggested flashcards depends on the
score: none at 85 or above, one from 60 to 84, and one or two below 60. The attempt is recorded
and the topic's mastery score is updated, while \code{/drill/save-cards} stores accepted cards.

\subsubsection{Functionality}
A wrong explanation becomes useful study material, since the flashcards produced target what
that student actually misunderstood, as summarised in Figure~\ref{fig:misconception}.

\shot{misconception}{Misconception detection and flashcard generation}{fig:misconception}

% ---- 7.6
\subsection{Flashcards and FSRS Review}
\subsubsection{Frontend}
The Flashcards page lists generated sets, each with its card and source counts, and a dialog
generates new sets from selected documents. The study view shows one question at a time and
links it to its source document and page. After revealing the answer, the student rates the
card as Again, Hard, Good or Easy. A Review page gathers all cards that are due.

\subsubsection{Backend}
Flashcards are generated from retrieved document chunks and validated before being saved.
Review sessions are handled through the \code{/api/v1/reviews} routes, with services to start a
session, fetch the next item, reveal an answer and rate an item. Scheduling sits behind an
abstract \code{ReviewScheduler} interface. Its default implementation uses the \code{fsrs}
library (FSRS version~6) with a desired retention of 90\% and a maximum interval of 36,500 days.

\subsubsection{Functionality}
FSRS models each card using its difficulty, stability and retrievability, as shown in
Figure~\ref{fig:fsrs}. Cards the student finds hard return sooner, and easy cards return later,
so revision time goes where it is most needed.

\shot[0.62\linewidth]{fsrs_cycle}{FSRS review cycle}{fig:fsrs}

\shot[0.72\linewidth]{flashcard}{Flashcard linked back to its source page}{fig:flashcard}

% ---- 7.7
\subsection{Knowledge Hub}
\subsubsection{Frontend}
The Knowledge Hub is a chat page with a list of saved conversations, suggested prompts and a
scope selector that limits the search to chosen documents or covers all of them. Each answer
shows its source references and offers study handoffs: Make Flashcard and Test Knowledge.

\subsubsection{Backend}
The question is processed by a retrieval service that scores document chunks using keyword
and vector similarity and then filters the results for diversity. In the current build, the
stored vectors are deterministic placeholder values rather than AI embeddings, so ranking
depends mainly on keyword matching. The selected chunks are passed to the model as untrusted
reference data, and the system prompt instructs the model not to follow any instructions found
inside them. A citation validator then checks each citation against the chunks that were
actually retrieved. Responses are streamed to the interface as server-sent events.

\subsubsection{Functionality}
Students can ask questions across their own material and check where each claim came from,
as shown in Figure~\ref{fig:hub}.

\shot{rag_flow}{Knowledge Hub answer flow}{fig:ragflow}

\shot{knowledge_hub}{Answer grounded in an uploaded chapter, with source marker [S1]}{fig:hub}

% ---- 7.8
\subsection{Quizzes}
\subsubsection{Frontend}
The Quizzes page lists generated quizzes with their difficulty, question count and sources.
The quiz detail page shows past attempts and an overview of the questions, and dedicated pages
handle taking a quiz and viewing its results.

\subsubsection{Backend}
The quiz service retrieves chunks from the selected documents, generates questions at the
requested difficulty and validates them before saving. Multiple-choice and true/false questions
are supported. The attempt service starts an attempt, records each answer and scores the attempt
on submission.

\subsubsection{Functionality}
Students can assess a whole lecture or document at any time and keep a record of attempts, which
also feeds into their progress statistics.

\shot[0.8\linewidth]{quiz}{Generated quiz questions}{fig:quiz}

% ---- 7.9
\subsection{Dashboard and Learning Progress}
\subsubsection{Frontend}
The Dashboard opens with the number of reviews due and a summary of the student's knowledge.
The Progress page reports review workload, daily activity, quiz accuracy and tracked knowledge.
Charts drawn with Recharts show reviews and quiz attempts over 7, 30 or 90 days, with a table view
as an alternative.

\subsubsection{Backend}
A learning analytics service provides the dashboard summary, review, quiz and flashcard statistics,
daily activity, and progress per document, per concept and per study session, through the
\code{/api/v1/learning} routes.

\subsubsection{Functionality}
The figures are computed from recorded activity, giving students a realistic view of what they have
studied and what needs attention next.

\shot{dashboard}{Student dashboard}{fig:dashboard}

\shot{progress}{Study consistency over time}{fig:progress}

% ---- 7.10
\subsection{AI Providers, Settings and Security}
\subsubsection{Frontend}
The Settings page is organised into tabs, including Account, Appearance, AI \& Providers and
Security. The AI \& Providers tab shows a card for each provider, where the student can save and test
a personal key and choose the active provider. Only a short hint of a saved key is displayed.

\subsubsection{Backend}
Personal API keys are encrypted with AES-256-GCM before being stored. Passwords are hashed using
PBKDF2-HMAC-SHA256 with 600,000 rounds, and authenticated requests to the application API use JSON
Web Tokens. Routes under \code{/api/v1/account/byok} save, delete and test credentials.

\subsubsection{Functionality}
Students can use their preferred cloud provider or a local Ollama model and can change between them
without affecting their study material.

\shot{providers}{Provider keys in the Settings page}{fig:providers}

% ------------------------------------------------------------------ 8
\section{Existing Related Work}
\label{sec:related}
Several existing tools address parts of the problem that RecaLL targets. A concise comparison based
on the vendors' published information is presented below. These products change quickly, so the
descriptions reflect their documentation at the time of writing.

\begin{enumerate}[leftmargin=*]
  \item \textbf{NotebookLM} (\url{https://notebooklm.google}) --- Google's AI research assistant
  answers from uploaded sources with inline citations, and can generate flashcards and practice quizzes.
  It is the closest tool in scope. Its published features emphasise working with sources, whereas RecaLL
  adds a scheduled revision path through FSRS and handwritten-style Cornell notes in a desktop application.
  \item \textbf{ChatPDF} (\url{https://www.chatpdf.com}) --- A chat tool that answers questions about
  uploaded PDFs, with clickable citations that lead to the relevant passage. It is convenient for finding
  information, but its focus is answering the student's questions rather than questioning the student.
  \item \textbf{Anki} (\url{https://apps.ankiweb.net}) --- A widely used open-source flashcard application,
  which has offered FSRS as an optional scheduler since version 23.10. Its scheduling is strong, but students
  generally create or import their own cards, whereas RecaLL generates cards from the source material.
  \item \textbf{Quizlet} (\url{https://quizlet.com}) --- A popular study platform whose AI tools generate
  study guides and flashcards from course materials, along with an adaptive Learn mode. It is a hosted
  service, whereas RecaLL keeps its data locally and lets the student choose the AI model.
  \item \textbf{RemNote} (\url{https://www.remnote.com}) --- A note-taking and spaced-repetition application
  that generates flashcards from notes and PDFs and supports FSRS. It is similar in its aim of connecting notes
  with revision. RecaLL differs in its Socratic drills, which turn a student's explanation into targeted
  flashcards, and in its option of a local model.
\end{enumerate}

Each of these tools is strong in its own area. RecaLL's contribution is to join the complete path, from
structuring a PDF to note-making, Socratic questioning, flashcard generation and FSRS revision, in one desktop
application with a choice of AI provider.

% ------------------------------------------------------------------ 9
\section{Conclusion}
RecaLL began as a response to a familiar problem: students read their textbooks carefully and still forget
what they read. The project developed into a desktop application that structures academic PDFs, writes Cornell
notes, questions the student through Socratic drills, turns gaps into flashcards and schedules revision with
FSRS. The Knowledge Hub and quizzes support this path, and the progress page keeps the student's activity visible.

The project also has clear limitations. Scanned PDFs require the optional Marker extractor, which is not part of
the default installation. Document search currently relies mainly on keyword matching, because AI embeddings are
not yet generated for stored chunks. Cloud AI requires an API key and an internet connection, local AI speed depends
on the student's hardware, and AI-generated content can still need checking against the source.

These limitations also indicate the direction of future work: built-in OCR for scanned documents, AI embeddings for
semantic search, smarter scheduling across subjects, and web and mobile access. Beyond the finished software, the
project reinforced the value of checking a system against its real behaviour and real documents, rather than relying
on what the code or its documentation appears to promise.

\end{document}
